\documentclass[../../../include/open-logic-section]{subfiles}

\begin{document}

\olfileid{cmp}{rec}{not}
\olsection{原始递归记号}

有了对原始递归函数的精确归纳描述，其优点之一是我们能够系统地描述它们。
例如，我们可以为每个这样的函数指定一个“记号”，具体如下。
使用符号 $\Zero$、$\Succ$ 和 $\Proj{n}{i}$ 分别表示零函数、后继函数和投影函数。
现在假设 $h$ 是由 $k$ 元函数~$f$ 和 $n$ 元函数 $g_0$，\dots，~$g_{k-1}$ 复合定义的，
并且我们已经为后者的这些函数指定了记号 $F$、$G_0$，\dots，~$G_{k-1}$。
那么，使用新符号 $\fn{Comp}_{k,n}$，我们可以用
$\fn{Comp}_{k,n}[F,G_0,\dots,G_{k-1}]$
来表示函数 $h$。
% Part: computability
% Chapter: recursive-functions
% Section: notations-pr-functions

对于由原始递归定义的函数，我们可以使用类似的记号。
假设 $(k+1)$ 元函数~$h$ 是由 $k$ 元函数~$f$ 和 $(k+2)$ 元函数~$g$ 经原始递归定义的，
并且分配给 $f$ 和~$g$ 的记号分别是 $F$ 和~$G$。
那么分配给~$h$ 的记号是 $\fn{Rec}_k[F,G]$。

回顾一下，加法函数由原始递归定义为
\begin{align*}
  \Add(x_0, 0) & = \Proj{1}{0}(x_0) = x_0\\
  \Add(x_0, y+1) & = \Succ(\Proj{3}{2}(x_0, y, \Add(x_0, y))) = \Add(x_0, y) +1
\end{align*}
这里，$f$ 对应于 $\Proj{1}{0}$，而 $g$ 对应于 $\Succ(\Proj{3}{2}(x_0, y, z))$，
它被指定的记号是 $\fn{Comp}_{1,3}[\Succ,\Proj{3}{2}]$，因为它是通过将 $1$ 元函数~$\Succ$ 和 $3$ 元函数~$\Proj{3}{2}$ 复合而定义的函数。
由此，我们可以用
\[
\fn{Rec}_1[\Proj{1}{0},\fn{Comp}_{1,3}[\Succ,\Proj{3}{2}]].
\]
来表示加法函数。
引入这些记号有时很有用，例如在枚举原始递归函数时。

\begin{prob}
  给出~$\Mult$ 的完整原始递归记号。
\end{prob}

\end{document}